\section{完整智能体栈：Perception, Cognition, Action}

上一节给出 autonomous agent 的广义范围，本节把这个概念拆成可调试的工程接口。读这套三层结构时要关注每层接收什么证据、输出什么状态，以及错误如何沿闭环传播；只有这样，机器人与 LLM Agent 的类比才不会停留在术语层。

Stone 在课里把完整智能体分成三个层次：Perception、Cognition、Action。这个拆法在机器人和 LLM 系统里都非常有用，因为它清晰地分离了 ``看见''、``想清楚''、``做出来'' 三个失败源。

\begin{figure}[H]
\centering
\includegraphics[width=0.95\textwidth]{diagrams/agent-loop.png}
\caption{讲义整理图：Perception、Cognition、Action 与环境反馈闭环}
\label{fig:agent-loop}
\end{figure}

\begin{knowledgebox}{读图：三层不是单向流水线}
Action 改变环境，新的结果再进入 Perception；队友、人类和对手都属于反馈的一部分。工程上应分别记录感知输入、内部 belief/plan 和执行结果，使失败能归因。若只保存最终文本，就无法判断是状态读错、计划错，还是动作接口执行错。
\end{knowledgebox}

\begin{table}[H]
\centering
\begin{tabular}{p{2.8cm}p{4.8cm}p{4.8cm}}
\toprule
\textbf{模块} & \textbf{在课程中的含义} & \textbf{在 LLM Agent 中的对应} \\
\midrule
Perception & 从 raw sensors 到可用状态表示 & context parsing, tool result interpretation \\
Cognition & 规划、推理、队友/对手建模、决策 & decomposition, planning, policy selection \\
Action & 从决策到可执行控制信号 & API call, code execution, UI action \\
\bottomrule
\end{tabular}
\caption{Perception-Cognition-Action 对应关系}
\end{table}

\begin{knowledgebox}{为什么要强制分层}
如果把系统全部扔进 end-to-end black box，会导致以下后果：
\begin{itemize}
    \item 出错时无法定位是感知问题还是决策问题；
    \item 奖励设计与评测指标很难有针对性；
    \item 安全约束无法插在合适位置（例如动作前的 rule check）。
\end{itemize}
\end{knowledgebox}

课程里还强调了 Cognition 层不只是 planning。它也包括 teammate/opponent modeling、coordination、tactical adaptation。这一点对多 Agent 任务尤其重要，因为策略质量不再只取决于自身状态，还取决于他人的策略分布。

\begin{importantbox}{智能体工程的一个实践准则}
在复杂任务里，先保证 Perception 和 Action 管道稳，再提高 Cognition 复杂度，通常比 ``一开始就上最复杂 planner'' 更容易收敛。
\end{importantbox}

\subsection{本章小结}
Perception-Cognition-Action 不是过时框架，而是今天构建可调试 Agent 的最低工程骨架。

